% !PRINTED 527
\weberpage{527}
Likewise 1 and 3 cannot adjoin one another cyclically, and there remain therefore only the following:
\weq{4}{(0,1,2,3), \quad (0,3,2,1), \quad (0,3)(1,2), \quad (0,1)(2,3).}
But if one of these permutations is added to (3), then all the others must also be admitted, for we have
\[
\begin{gathered}
(0,2)(0,1,2,3) = (0,3)(1,2) \\
(1,3)(0,1,2,3) = (0,1)(2,3) \\
(0,2)(1,3)(0,1,2,3) = (0,3,2,1),
\end{gathered}
\]
and one can therefore, by means of the transpositions $(0,2)$ and $(1,3)$, derive each of the permutations (4) from each of the others. From the assumption, then, that a transposition occurs at all in a transitive group different from $P$ and $Q$, it follows that this must be the group of the eighth degree
\weq{5}{\begin{gathered}
P_1 = 1, \quad (0,2), \quad (1,3), \quad (0,2)(1,3), \quad (0,1)(2,3), \quad (0,3)(1,2) \\
(0,1,2,3), \quad (0,3,2,1)
\end{gathered}}
or at least a group conjugate with it.

That $P_1$ really is a group one easily recognises by composing any two of its permutations.

There are three distinct groups conjugate with $P_1$, which one obtains by starting, instead of from $(0,2)$, $(1,3)$, from the transpositions $(0,1)$, $(2,3)$ or $(0,3)$, $(1,2)$. The intersection of these three groups is a normal divisor of $P$ and at the same time of $Q$, namely
\weq{6}{Q_1 = 1, \quad (0,2)(1,3), \quad (0,1)(2,3), \quad (0,3)(1,2).}
If a transitive group different from $P$ and $Q$ contains no transposition, then it either contains no four-term cycle and is then identical with (6), or it contains a four-term cycle and then contains also the whole cyclic group, such as
\weq{7}{P_2 = 1, \quad (0,1,2,3), \quad (0,2)(1,3), \quad (0,3,2,1).}
If such a group contains yet a second four-term cycle, say $(0,1,3,2)$, then on account of $(0,1,2,3)(0,1,3,2) = (0,3,1)$ it is the symmetric group. But if, besides $P_2$, it contains also one of the permutations of (6), say $(0,1)(2,3)$, then on account of
\[
(0,1,2,3)(0,1)(2,3) = (1,3)
\]
it is identical with $P_1$.
